\section{Ablation I: Contribution of CAT and MIL Terms}
\label{appendix:catmil-components}

Unless stated otherwise, all ablation experiments use the same nnU-Net
configuration, data splits, preprocessing, and training schedule as the
main experiments. To isolate the effect of each component, three variants
are compared: (i) CAT, which adds only the structure-level
component-adaptive term to the base loss, (ii) MIL, which adds only the
lesion-level supervision term, and (iii) CATMIL, which combines both terms.

\Cref{tab:ablation_cat_mil} reports the mean performance of the five fold
models on the MSLesSeg test set, and the differences below are given with
95\% patient-level CIs as in \cref{tab:paired_differences}.

\textbf{The MIL term accounts for the detection gain.} MIL alone raises
small-lesion recall from 0.7956 (DiceCE) to 0.8712 (+0.076 [+0.033,
+0.148]) and reduces missed lesions per case by 2.42 [0.87, 5.26], in all
six patients. These gains are as large as those of CATMIL, and CATMIL and
MIL do not differ on any metric: small-lesion recall differs by +0.002
[$-0.012$, +0.019], lesion-wise F1 by +0.008 [$-0.020$, +0.031], DSC by
+0.003 [$-0.016$, +0.021], and HD95 by $-1.26$~mm [$-6.82$, +1.55]. The
lower lesion-wise F1 of CATMIL is also already present with MIL alone
(0.7492).

\textbf{The CAT term alone changes lesion detection little.} CAT alone
reaches a small-lesion recall of 0.8230, a difference from DiceCE of
+0.027 [$-0.010$, +0.081] whose CI includes zero, and it keeps the
lesion-wise F1 of DiceCE (0.8386 against 0.8433). Its effect on detection
is thus close to that of the Tversky-based losses, consistent with a term
that rebalances voxel weights between lesions but, unlike the MIL term,
does not require each lesion to be detected.

The mean DSC and HD95 of CATMIL are the best of the three variants, and
CATMIL has the lowest false-positive blob fraction, but these differences
lie within the variation between patients and fold models. On this test
set, the ablation therefore supports the MIL term as the source of the
detection gain, and it does not show a measurable benefit of adding the
CAT term to it. The CAT term is kept in CATMIL because of its design
rationale, size-balanced voxel supervision within each lesion
(\cref{appendix:catmil-properties}), and because the loss-weight analysis
below shows that its weight shifts the balance between voxel-level accuracy
and lesion-wise precision; whether it contributes measurably on other data
remains to be tested.

\begin{table}[H]
\centering
\caption{Ablation study on the individual and combined contribution of the
CAT and MIL terms on the MSLesSeg test set. FP Vol: false-positive volume
(mm$^3$); FP Blob: false-positive blob fraction.}
\label{tab:ablation_cat_mil}
\begin{tabular}{lcccccc}
\toprule
Model & DSC & HD95 $\downarrow$ & Lesion F1 & Small Lesion Recall & FP Vol $\downarrow$ & FP Blob $\downarrow$ \\
\midrule
CAT    & 0.7812 & 8.6844 & \textbf{0.8386} & 0.8230 & \textbf{1528.72} & 0.2277 \\
MIL    & 0.7805 & 9.2375 & 0.7492 & 0.8712 & 1539.77 & 0.2195 \\
CATMIL & \textbf{0.7834} & \textbf{7.9817} & 0.7571 & \textbf{0.8730} & 1536.88 & \textbf{0.2094} \\
\bottomrule
\end{tabular}
\end{table}

\section{Ablation II: Sensitivity to Loss Weights}
\label{appendix:catmil-loss-weights}

This section analyzes the impact of the weighting coefficients
$\lambda_{\text{CAT}}$ and $\lambda_{\text{MIL}}$. The configurations are
named CATMIL$xy$, where $\lambda_{\text{CAT}}=0.x$ and
$\lambda_{\text{MIL}}=0.y$; CATMIL0102 is the configuration used in the
main experiments. Its weights were fixed before this analysis was run, and
the analysis was not used to select them: it is a post hoc check of how
sensitive the results are to the choice of weights. Unlike in the other
experiments, each configuration in this analysis was trained once on the
complete development set, without cross-validation, and evaluated on the
same test set. The values of CATMIL0102 therefore differ from those of
CATMIL in \cref{tab:loss_comparison_small_lesion,tab:ablation_cat_mil},
which are means over the five cross-validation models.

\begin{table}[H]
\centering
\caption{Sensitivity of the CATMIL objective to $\lambda_{\text{CAT}}$ and
$\lambda_{\text{MIL}}$ on the MSLesSeg test set.}
\label{tab:ablation_lambda}
\begin{tabular}{lccccc}
\toprule
Model & DSC & HD95 $\downarrow$ & Lesion F1 & Small Recall & FP Vol $\downarrow$ \\
\midrule
CATMIL0101 & 0.7675 & 9.3881 & \underline{0.7720} & 0.8514 & 2292.17 \\
CATMIL0102 & \textbf{0.7954} & \textbf{5.9771} & 0.7452 & \underline{0.8711} & \underline{1345.83} \\
CATMIL0201 & \underline{0.7921} & \underline{6.7964} & \textbf{0.7754} & \textbf{0.8784} & \textbf{1233.92} \\
CATMIL0202 & 0.7663 & 9.4507 & 0.7632 & 0.8699 & 1729.42 \\
\bottomrule
\end{tabular}
\end{table}

The configuration used in the main experiments (CATMIL0102), in which the
MIL term has the larger weight, gives the best global segmentation
performance, with the highest DSC (0.7954) and lowest HD95 (5.9771~mm), but
also the lowest lesion-wise F1 (0.7452). Exchanging the two weights
(CATMIL0201) improves lesion-level behavior, with the highest lesion F1
(0.7754), the highest small lesion recall (0.8784), and the lowest FP
volume (1233.92~mm$^3$), at the cost of slightly lower DSC and boundary
accuracy. The two configurations with equal weights (CATMIL0101 and
CATMIL0202) are clearly worse in DSC, HD95, and FP volume. Small lesion
recall stays between 0.85 and 0.88 in all four configurations, so the
detection gain does not depend on a precise choice of the weights in this
range; the ratio of $\lambda_{\text{CAT}}$ to $\lambda_{\text{MIL}}$ mainly
shifts the balance between voxel-level accuracy and lesion-wise precision.

\section{Lesion-Profile Characteristics of Test Cases}
\label{appendix:test-case-lesion-profile}

Each MSLesSeg test case is summarized using the ground-truth lesion mask
(\cref{tab:test_case_lesion_profile}). Lesions are counted by applying
connected-component analysis with 6-connectivity to the binary lesion
mask, the same definition of a lesion as in the proposed loss and in the
evaluation metrics; all components are counted regardless of their size.
Lesion volume is the number of voxels inside each lesion multiplied by the
voxel volume; at the 1~mm isotropic spacing of MSLesSeg, one voxel
corresponds to 1~mm$^3$. The total volume is the sum of all lesion volumes
in one case, and the median volume represents the typical lesion size in
that case.

Lesion contrast is computed locally on the FLAIR image after z-score
normalization within the brain. For each lesion, a surrounding background
ring is created within 3~mm around the lesion, excluding the lesion itself.
The contrast value is the absolute difference between the median intensity
inside the lesion and the median intensity in this ring, divided by the
median absolute deviation of the ring intensity. A higher value indicates
that the lesion is more distinguishable from its local background.

Boundary sharpness is computed from the image gradient around the lesion
boundary. The normalized FLAIR image is smoothed with a Gaussian filter of
1~mm, and the sharpness is the median gradient magnitude in a band of one
voxel on each side of the lesion boundary. A higher value indicates a
clearer lesion boundary.

Contrast and sharpness are not estimated reliably for the smallest lesions
and are therefore computed only for lesions of at least 10 voxels. The
contrast and sharpness in \cref{tab:test_case_lesion_profile} are means
over these lesions, whereas the lesion count and the volume statistics
include all lesions.

\begin{table}[H]
\centering
\caption{Lesion-profile characteristics of the MSLesSeg test cases. N:
number of lesions; $\leq$50 and $\leq$100: percentage of lesions of at most
50 and 100~mm$^3$.}
\label{tab:test_case_lesion_profile}
\begin{tabular}{lrrrrrrr}
\toprule
\textbf{Case} & \textbf{N} & \textbf{Total vol.} & \textbf{Median vol.}
& \textbf{$\leq$50} & \textbf{$\leq$100} & \textbf{Contrast} & \textbf{Sharpness} \\
& & \textbf{(mm$^3$)} & \textbf{(mm$^3$)} & \textbf{(\%)} & \textbf{(\%)} & & \\
\midrule
P13\_T1 & 39 & 9677 & 83.0 & 10.3 & 53.8 & 4.11 & 8.41 \\
P13\_T2 & 61 & 10268 & 68.0 & 45.9 & 65.6 & 3.85 & 8.34 \\
P26\_T1 & 43 & 6639 & 106.0 & 20.9 & 46.5 & 3.77 & 9.16 \\
P2\_T1  & 16 & 2261 & 92.5 & 31.2 & 50.0 & 4.56 & 9.31 \\
P2\_T2  & 20 & 3286 & 87.5 & 20.0 & 55.0 & 3.76 & 10.30 \\
P2\_T3  & 14 & 2623 & 134.0 & 14.3 & 28.6 & 4.61 & 11.68 \\
P2\_T4  & 18 & 2940 & 87.0 & 16.7 & 55.6 & 4.45 & 10.86 \\
P47\_T1 & 59 & 18736 & 67.0 & 42.4 & 71.2 & 4.26 & 9.92 \\
P49\_T1 & 30 & 34709 & 83.5 & 30.0 & 56.7 & 3.65 & 8.55 \\
P49\_T2 & 19 & 7763 & 173.0 & 0.0 & 26.3 & 3.34 & 7.38 \\
P52\_T1 & 28 & 3306 & 70.5 & 42.9 & 67.9 & 5.24 & 11.02 \\
P52\_T2 & 14 & 2843 & 147.0 & 7.1 & 42.9 & 4.04 & 6.91 \\
\bottomrule
\end{tabular}
\end{table}

\section{Additional Lesion-Level Analyses}
\label{appendix:lesion-level-analyses}

This section reports lesion-level analyses that complement the size
analysis in the main text. All analyses use the same lesion definition as
the evaluation metrics: 6-connected components of the reference mask, with
a lesion counted as detected if any predicted voxel overlaps it. Values are
means over the five fold models.

\paragraph{Operating points under component filtering.}
PP5 uses a single filter size. \Cref{fig:operating_curve} shows how lesion
recall and the number of false-positive components per case change when
the filter size is varied. In the left panel, predicted components smaller
than $s$ voxels are removed, with $s$ from 0 to 20; in the right panel, the
probability threshold $t$ is varied from 0.1 to 0.9 without filtering. A
curve that lies higher and further to the left detects more lesions for
the same number of false alarms.

The CATMIL curve lies above the baseline curves across the whole range of
filter sizes. Without filtering, CATMIL produces more false-positive
components than the baselines (20.3 per case, against 5.4--7.0), which is
the cause of its lower Lesion F1. These components are small and are
removed at little cost in recall. With $s = 20$, CATMIL detects 84.5\% of
lesions at 3.3 false-positive components per case, more lesions with fewer
false alarms than unfiltered DiceCE (83.5\% at 5.4). For any false-alarm
budget of at least 3.3 components per case, a filtered CATMIL model
therefore detects more lesions than any of the baselines. Raising the
probability threshold is less effective than a size filter, because most
of the false-positive components are confident predictions rather than
borderline ones. The margin over the Tversky-based losses at a fixed
false-alarm level is small, however, and is not supported by the
patient-level analysis at matched precision (\cref{tab:matched_precision}).

\paragraph{Joint variation of threshold and filter size.}
\Cref{fig:operating_grid} varies the probability threshold and the minimum
component size jointly and plots lesion recall against lesion-wise
precision, the two quantities that make up lesion-wise F1. Each faint point
is one setting; the lines connect the settings that are not outperformed
in both precision and recall by another setting of the same loss. On
MSLesSeg, the CATMIL frontier lies above that of DiceCE over most of the
precision range and close to those of the Tversky-based losses, which is
the pattern quantified in \cref{tab:matched_precision}. On 3D-MR-MS, the
default output of CATMIL (filled circle) lies far to the left of its
frontier: its precision of 0.37 is due to the many small components, and
the frontier reaches the precision of the baselines with similar recall.
In this analysis, the minimum component size is given in voxels on both
datasets. Expressing it in mm$^3$ instead, that is, removing components of
the same physical volume on both datasets, transfers poorly: because the
voxels of 3D-MR-MS are 5.7 times smaller, a filter of 15~mm$^3$ removes
85 voxels and with them most of the small lesions. The false-positive
components of CATMIL are therefore small in voxels rather than in
millimeters, and the filter size should be chosen for each image
resolution.

\begin{figure}[htbp]
    \centering
    \includegraphics[width=\textwidth]{figures/operating_grid.pdf}
    \caption{Lesion recall against lesion-wise precision for all
    combinations of probability threshold ($t = 0.1, \dots, 0.9$) and
    minimum component size ($s = 1$ to 170 voxels). Top: small-lesion
    recall ($\leq 150$ voxels); bottom: recall of all lesions. Faint
    points: individual settings; lines: settings not outperformed in both
    precision and recall by another setting of the same loss. Filled
    markers: default output ($t = 0.5$, $s = 1$); open markers: PP5
    ($t = 0.5$, $s = 5$), shown for CATMIL and DiceCE. Metrics are computed
    per case and averaged over cases and the five fold models.}
    \label{fig:operating_grid}
\end{figure}

\begin{figure}[htbp]
    \centering
    \includegraphics[width=\textwidth]{figures/operating_curve.pdf}
    \caption{Lesion recall against false-positive components per case on
    the MSLesSeg test set. Left: predicted components smaller than $s$
    voxels are removed ($s = 0, 2, 3, 5, 8, 10, 15, 20$) at probability
    threshold 0.5. Right: the probability threshold $t$ is varied from 0.1
    to 0.9 without filtering. Open circles: unfiltered predictions at
    $t = 0.5$; stars: PP5 ($s = 5$). A false-positive component is a
    predicted connected component that does not overlap any reference
    lesion. Recall is pooled over the test lesions; values are means over
    the five fold models.}
    \label{fig:operating_curve}
\end{figure}

\paragraph{Lesion-level analyses on the 3D-MR-MS dataset.}
\Cref{fig:recall_vs_size_curve_3dmrms,fig:lesion_fate_3dmrms,fig:operating_curve_3dmrms,fig:max_prob_by_size_3dmrms}
repeat the analyses of
\cref{fig:recall_vs_size_curve,fig:lesion_fate,fig:operating_curve,fig:max_prob_by_size}
on the 3D-MR-MS dataset, and are read in the same way. Because the voxels
of this dataset are anisotropic ($0.47 \times 0.47 \times 0.8$~mm), lesion
sizes are given in mm$^3$ rather than in voxels.

The recall curve (\cref{fig:recall_vs_size_curve_3dmrms}) reproduces the
main result. CATMIL detects 51\% of lesions of at most 10~mm$^3$, compared
with 40--41\% for the baselines, and 93\% of lesions of 11--50~mm$^3$,
compared with 86--88\%. Small lesions are far more frequent in this
dataset: 182 of the 389 reference lesions are at most 10~mm$^3$, so the
gain in this range accounts for most of the overall increase in lesion
recall, from 0.66--0.67 to 0.74. Above 100~mm$^3$, CATMIL detects slightly
fewer lesions than the baselines (95\% against 98--99\%), a difference of
about one of the 34 lesions in this range.

The paired comparison (\cref{fig:lesion_fate_3dmrms}) is again one-sided.
CATMIL detects 33.2 lesions per fold model that DiceCE misses, against 2.4
in the opposite direction, and 30.4 against 3.8 for each of the
Tversky-based losses. About 97\% of the lesions recovered by CATMIL in each
comparison are at most 50~mm$^3$.

The operating curve (\cref{fig:operating_curve_3dmrms}) shows that the
component-level cost of CATMIL is larger on this dataset. Without
filtering, CATMIL produces 150 false-positive components per case, compared
with 14--19 for the baselines, which explains its low Lesion F1 in
\cref{tab:external_eval_3dmrms}. The size filter again removes most of
them: PP5 reduces the count to 27 per case with a recall of 0.70, and a
filter of 20 voxels reduces it to 9.7 per case with a recall of 0.657. At
this setting, CATMIL has about as many false-positive components as DiceCE
after PP5 (9.5 per case) and a higher recall (0.648 for DiceCE). The filter
size required to control false positives is therefore dataset dependent,
and a larger filter is needed at the finer in-plane resolution of
3D-MR-MS.

The probability levels (\cref{fig:max_prob_by_size_3dmrms}) show the same
mechanism as on MSLesSeg. For lesions of at most 10~mm$^3$, the fraction
with no response falls from 57--58\% for the baselines to 39\% for CATMIL.

\begin{figure}[htbp]
    \centering
    \includegraphics[width=0.85\textwidth]{figures/recall_vs_size_curve_3dmrms.pdf}
    \caption{Lesion recall against lesion size on the 3D-MR-MS test set.
    Layout as in \cref{fig:recall_vs_size_curve}.}
    \label{fig:recall_vs_size_curve_3dmrms}
\end{figure}

\begin{figure}[htbp]
    \centering
    \includegraphics[width=\textwidth]{figures/lesion_fate_3dmrms.pdf}
    \caption{Lesions detected by only one of two losses on the 3D-MR-MS
    test set, by lesion size. Layout as in \cref{fig:lesion_fate}.}
    \label{fig:lesion_fate_3dmrms}
\end{figure}

\begin{figure}[htbp]
    \centering
    \includegraphics[width=\textwidth]{figures/operating_curve_3dmrms.pdf}
    \caption{Lesion recall against false-positive components per case on
    the 3D-MR-MS test set. Layout as in \cref{fig:operating_curve}.}
    \label{fig:operating_curve_3dmrms}
\end{figure}

\begin{figure}[htbp]
    \centering
    \includegraphics[width=0.75\textwidth]{figures/max_prob_by_size_3dmrms.pdf}
    \caption{Maximum foreground probability inside each reference lesion on
    the 3D-MR-MS test set, by lesion size. Layout as in
    \cref{fig:max_prob_by_size}.}
    \label{fig:max_prob_by_size_3dmrms}
\end{figure}

\section{Gradient and Boundedness Properties of the CATMIL Objective}
\label{appendix:catmil-properties}

This section derives the properties of the CAT and MIL terms stated in
\cref{sec:method}, using the notation of that section. The components $C_k$
and the weights $w_i$ are computed from the ground-truth mask and are
constants with respect to the predictions $p_i$.

\paragraph{Gradient of the CAT term.}
Let
\begin{equation}
A = \mathrm{TP}_{w}+\delta,
\qquad
B = A + \alpha\, \mathrm{FP}_{w} + \beta\, \mathrm{FN}_{w},
\end{equation}
so that $\mathcal{L}_{\mathrm{CAT}} = 1 - A/B$. For a foreground voxel
$j \in C_k$, $\partial \mathrm{TP}_{w}/\partial p_j = w_j$,
$\partial \mathrm{FP}_{w}/\partial p_j = 0$, and
$\partial \mathrm{FN}_{w}/\partial p_j = -w_j$, hence
$\partial A/\partial p_j = w_j$ and
$\partial B/\partial p_j = (1-\beta)\, w_j$. The quotient rule gives
\begin{equation}
\frac{\partial \mathcal{L}_{\mathrm{CAT}}}{\partial p_j}
=
-\frac{w_j B - (1-\beta)\, w_j A}{B^2}
=
-c\, w_j ,
\qquad
c
=
\frac{\beta A + \alpha\, \mathrm{FP}_{w} + \beta\, \mathrm{FN}_{w}}{B^2}
> 0 .
\end{equation}
The factor $c$ is the same for all foreground voxels, so the gradient at a
lesion voxel is scaled only by its weight
$w_j = (|C_k|+\varepsilon_{c})^{-\gamma}$. Summing over the component, the
total gradient magnitude received by lesion $C_k$ is
\begin{equation}
\sum_{j \in C_k}
\left| \frac{\partial \mathcal{L}_{\mathrm{CAT}}}{\partial p_j} \right|
=
c\, |C_k|\, (|C_k|+\varepsilon_{c})^{-\gamma},
\end{equation}
which is proportional to $|C_k|$ for $\gamma=0$ and equal to
$c\,\rho_k < c$ for $\gamma=1$, with $\rho_k$ the lesion weight of
\cref{eq:sl_cat_lesion_weight}.

For a background voxel $j$ with $g_j=0$, only $\mathrm{FP}_{w}$ depends on
$p_j$, so $\partial A/\partial p_j = 0$ and
$\partial B/\partial p_j = \alpha\, w_{\mathrm{bg}}$, and
\begin{equation}
\frac{\partial \mathcal{L}_{\mathrm{CAT}}}{\partial p_j}
=
\frac{\alpha\, w_{\mathrm{bg}}\, A}{B^2}
\geq 0 .
\end{equation}
Increasing the predicted probability at a background voxel therefore
increases the loss.

\paragraph{Gradient of the MIL term.}
With the clipping of \cref{eq:sl_mil_loss}, the contribution of lesion
$C_k$ is $-\log\bigl(\max(s_k,\varepsilon)\bigr)$ and remains finite. Let
$m_k = \arg\max_{i \in C_k} p_i$ be the maximal voxel of lesion $C_k$.
Away from ties,
\begin{equation}
\frac{\partial \mathcal{L}_{\mathrm{MIL}}}{\partial p_j}
=
\begin{cases}
-\dfrac{1}{K s_k},
& \text{if } j=m_k \text{ for some } k \text{ and } s_k>\varepsilon, \\[2mm]
0, & \text{otherwise.}
\end{cases}
\end{equation}
The gradient is concentrated on one voxel per lesion, its magnitude does
not depend on $|C_k|$, and it grows as the strongest response $s_k$
decreases. When several voxels share the maximum, the hard maximum is not
differentiable but remains subdifferentiable, and automatic differentiation
assigns the gradient to a maximal voxel, as in max-pooling.

\paragraph{Boundedness and continuity.}
Since $\mathrm{TP}_{w}$, $\mathrm{FP}_{w}$, $\mathrm{FN}_{w}$, $\alpha$, and
$\beta$ are non-negative and $\delta>0$, $0 < A \leq B$, and therefore
$0 \leq \mathcal{L}_{\mathrm{CAT}} < 1$. Since
$\varepsilon \leq \max(s_k,\varepsilon) \leq 1$,
$0 \leq \mathcal{L}_{\mathrm{MIL}} \leq -\log \varepsilon$. For
non-negative coefficients the auxiliary part of the objective thus
satisfies
\begin{equation}
0
\leq
\lambda_{\mathrm{CAT}}\, \mathcal{L}_{\mathrm{CAT}}
+
\lambda_{\mathrm{MIL}}\, \mathcal{L}_{\mathrm{MIL}}
\leq
\lambda_{\mathrm{CAT}} - \lambda_{\mathrm{MIL}} \log \varepsilon .
\end{equation}
Both terms are continuous in the predictions. The CAT term is a ratio of
linear functions with a strictly positive denominator and is differentiable
everywhere. The MIL term is differentiable except at ties of the maximum
and at the clipping point $s_k=\varepsilon$, which is the same piecewise
behavior as that of max-pooling and does not prevent gradient-based
training.
